% Appendix C — Robustness Specifications
% Phase 10 cleanup: dropped stale v5.1 BelowIG three-tier subsection (binary scheme is the actual v6 design);
% updated v5.1 UncAnsCEO IV references to v6 three-component decomp; updated stale table refs to _decomp.

\section{Robustness Specifications}
\label{app:robustness}

This appendix collects auxiliary specifications and design-choice disclosures supporting the main \S\ref{sec:main} hypotheses HC and HFC.

\subsection{Binary Constraint Moderator: Rated-versus-Unrated}
\label{app:robustness:rated}

The HFC specification of \S\ref{sec:main:hfc} adopts the binary rated-versus-unrated partition of \citeA{faulkender2006} directly: the Unrated indicator equals one for firms without a Compustat credit rating, and the rated reference category pools investment-grade and below-investment-grade firms. The binary scheme follows \citeauthor{faulkender2006} verbatim and is consistent with the macro-shock asymmetry framing of \citeA{almeida2004}~\S II.D. We do not estimate a three-tier IG/BelowIG/Unrated specification: the main hypothesis localizes the precautionary cash response to the credit-constrained Unrated cell, and the rated half of the population is heterogeneous in ways that the binary scheme treats as nuisance variation absorbed into the rated-pool reference.

\subsection{Alternative Cash Dependent-Variable Specification (OPSW Log-Form)}
\label{app:robustness:opsw}

The body adopts the linear cash-to-assets dependent variable (\texttt{cheq/atq}) of \citeA{bates2009} as the standard precautionary-cash form. \citeA{opler1999} also use a log-form dependent variable, defined as $\log(\text{cheq}/\text{atq})$ (or, in some specifications, $\log(\text{cheq}/(\text{atq}-\text{cheq}))$ to avoid log-of-zero issues at the lower tail). The log form down-weights the upper tail of the cash distribution relative to the linear form and tests whether the precautionary association is driven by high-cash firms.

The present draft does not estimate the OPSW log-DV robustness; it is flagged in \S\ref{sec:conclusion} as a planned extension. The expectation, given the magnitude of the linear-form coefficients in \S\ref{sec:main:hc} (\texttt{UncResCEO} significant 12 of 12 at $p<0.10$), is that the qualitative result survives under the log specification, with attenuated point-estimate magnitudes due to the tail-compression. Future work will report the log-DV coefficients directly to confirm this expectation.

\subsection{Fixed-Effects Ladder Robustness}
\label{app:robustness:fe}

The four-step fixed-effects ladder of \S\ref{sec:framework:design} (Industry$+$Year, Firm$+$Year, Industry$+$Year-Quarter, Firm$+$Year-Quarter) is the within-regressor robustness frame applied throughout the body. Every column of every body table reports a different point on this ladder, so the FE-ladder robustness check is integrated rather than a separate appendix specification. The convergence pattern across the ladder---attenuation from industry to firm fixed effects, additional attenuation from year to year-quarter fixed effects---is the standard within-firm panel-OLS attenuation pattern (within-firm variation is a strict subset of cross-firm variation, so coefficients on a regressor with substantial across-firm dispersion shrink as fixed effects absorb that dispersion). The signs and significance of \texttt{UncResCEO} are robust to all four ladder steps in HC, HFC, and the cash-flow-volatility moderator; this robustness is documented in the body and not separately tabulated here.

\subsection{Full and QtrExp Measurement Methods}
\label{app:robustness:methods}

The body reports each of the HC, HFC, and CFvol suites under two parallel measurement methods. The Full method estimates the underlying CEO fixed-effect regression on the entire firm-quarter panel and recovers the negative of the CEO fixed effect as \texttt{ClarityCEO} and the regression residual as \texttt{UncResCEO} (the original \citeA{dzielinski2021} construction). The QtrExp method estimates an expanding-window CEO fixed effect on within-CEO data only, using only the calls of a given CEO available up to and including the focal call. The QtrExp method is strictly look-ahead-free: the recovered components for a focal call use no information from calls that follow it. The two methods agree in direction on all three speech components in all body suites and yield slightly stronger statistical evidence for the persistent-style \texttt{ClarityCEO} component in the QtrExp method (which uses fewer observations and a more demanding within-CEO identification window). The agreement between methods supports interpreting the recovered components as CEO-trait and call-state signals rather than as panel-fitted artifacts.
